\section{Gabarito e resoluções comentadas — Engenharia de Dados}

\begin{solucao}{11.01}\textbf{Gabarito: A.} A chave primária identifica unicamente cada tupla e não admite duplicidade. Ela não obriga ordenação física, não substitui todos os índices e não elimina integridade referencial.\end{solucao}
\begin{solucao}{11.02}\textbf{Gabarito: A.} Chave estrangeira representa referência a chave candidata/primária de outra relação e ajuda a garantir integridade referencial. Não é mecanismo de criptografia, backup ou paginação.\end{solucao}
\begin{solucao}{11.03}\textbf{Gabarito: A.} A 1FN trata atomicidade dos valores no contexto relacional e ausência de grupos repetitivos. Não elimina toda redundância nem exige BCNF, chave simples ou surrogate key.\end{solucao}
\begin{solucao}{11.04}\textbf{Gabarito: A.} \texttt{CREATE TABLE} define estrutura e pertence à DDL. SELECT consulta dados; INSERT, UPDATE e DELETE manipulam dados.\end{solucao}
\begin{solucao}{11.05}\textbf{Gabarito: A.} Tabela fato registra eventos e medidas na granularidade escolhida. Dimensões descrevem contexto. Índice, catálogo e view materializada não exercem esse papel por definição.\end{solucao}
\begin{solucao}{11.06}\textbf{Gabarito: A.} Bancos de documentos armazenam documentos semiestruturados, frequentemente JSON/BSON. Tabela 3FN é relacional; CSV não é SGBD; RAID e DNS não são bancos NoSQL.\end{solucao}
\begin{solucao}{11.07}\textbf{Gabarito: A.} ETL = Extract, Transform, Load. As demais expansões são fictícias.\end{solucao}
\begin{solucao}{11.08}\textbf{Gabarito: A.} Linhagem registra de onde o dado veio, quais transformações sofreu e para onde foi. Isso sustenta rastreabilidade e reprodutibilidade. As demais opções são metadados técnicos de outras áreas.\end{solucao}
\begin{solucao}{11.09}\textbf{Gabarito: A.} Durabilidade significa que, após commit, os efeitos persistem conforme o modelo de falhas do SGBD. Atomicidade, consistência e isolamento tratam outras propriedades; idempotência não integra ACID.\end{solucao}
\begin{solucao}{11.10}\textbf{Gabarito: A.} HAVING filtra grupos após agregação; WHERE filtra linhas antes da agregação. HAVING não cria tabela nem define chaves.\end{solucao}
\begin{solucao}{11.11}\textbf{Gabarito: A.} $C$ depende apenas de parte da chave composta, logo há dependência parcial, problema tratado pela 2FN. Não é, por isso só, dependência multivalorada ou de junção.\end{solucao}
\begin{solucao}{11.12}\textbf{Gabarito: A.} Se PedidoID determina ClienteID e ClienteID determina NomeCliente, há caminho transitivo de chave para atributo não-chave, típico problema de 3FN. Não é mero problema de 1FN.\end{solucao}
\begin{solucao}{11.13}\textbf{Gabarito: A.} LEFT JOIN preserva todas as linhas da relação à esquerda. INNER JOIN descartaria órgãos sem correspondência; CROSS JOIN produz produto cartesiano.\end{solucao}
\begin{solucao}{11.14}\textbf{Gabarito: A.} B-tree atende bem igualdade e intervalos em muitos SGBDs. Hash é mais associado a igualdade e não é universal; as demais alternativas não fazem sentido como índice relacional geral.\end{solucao}
\begin{solucao}{11.15}\textbf{Gabarito: A.} Baixa cardinalidade pode gerar baixa seletividade, fazendo o otimizador preferir scan conforme custo. Índices não são benefício automático e também têm custo de manutenção.\end{solucao}
\begin{solucao}{11.16}\textbf{Gabarito: A.} A granularidade define o significado de cada linha da fato. Sem isso, medidas podem ser somadas de modo incorreto e produzir dupla contagem.\end{solucao}
\begin{solucao}{11.17}\textbf{Gabarito: A.} Saldo é normalmente semi-aditivo: pode ser somado em algumas dimensões, mas não ao longo do tempo sem regra apropriada.\end{solucao}
\begin{solucao}{11.18}\textbf{Gabarito: A.} Em ELT, dados são carregados antes e transformados no destino. Extração continua existindo e transformação não é proibida.\end{solucao}
\begin{solucao}{11.19}\textbf{Gabarito: A.} Idempotência permite repetir processamento sem duplicar o efeito lógico. É central em pipelines com retry. As demais opções não resolvem duplicidade de efeito.\end{solucao}
\begin{solucao}{11.20}\textbf{Gabarito: A.} Completude verifica presença de dados esperados/obrigatórios. Unicidade, atualidade e consistência são dimensões distintas.\end{solucao}
\begin{solucao}{11.21}\textbf{Gabarito: A.} Atomicidade exige tudo ou nada para a transação. Isolamento trata concorrência e durabilidade trata persistência após commit.\end{solucao}
\begin{solucao}{11.22}\textbf{Gabarito: A.} Dirty read é leitura de dado ainda não confirmado. Não implica deadlock nem full scan.\end{solucao}
\begin{solucao}{11.23}\textbf{Gabarito: A.} O primeiro passo é verificar plano, estatísticas, seletividade e índice adequado. Criar índices indiscriminadamente pode piorar escrita e desperdiçar espaço.\end{solucao}
\begin{solucao}{11.24}\textbf{Gabarito: A.} Catálogo torna datasets encontráveis e compreensíveis, registrando esquema, significado, responsável e outras informações. Ele não substitui backup, segurança ou qualidade.\end{solucao}
\begin{solucao}{11.25}\textbf{Gabarito: A.} Contagens, hashes e totais de controle permitem reconciliar origem e destino e detectar perda/corrupção. Nome do arquivo sozinho é evidência fraca.\end{solucao}
\begin{solucao}{11.26}\textbf{Gabarito: A.} NomeAluno deve ser separado em relação de aluno para remover dependência parcial. Duplicar o nome aumenta anomalias; índice não corrige desenho lógico.\end{solucao}
\begin{solucao}{11.27}\textbf{Gabarito: A.} Ao exigir \texttt{t2.status='ATIVO'} no WHERE, linhas sem correspondência têm NULL e são eliminadas. Isso neutraliza a preservação típica do LEFT JOIN para esse predicado.\end{solucao}
\begin{solucao}{11.28}\textbf{Gabarito: A.} Lost update é anomalia de concorrência. Deve-se avaliar isolamento, locking, MVCC, versionamento otimista e padrão da atualização.\end{solucao}
\begin{solucao}{11.29}\textbf{Gabarito: A.} MVCC mantém versões para permitir leituras consistentes e reduzir bloqueios conforme SGBD. Não substitui backup/WAL nem elimina manutenção de versões antigas.\end{solucao}
\begin{solucao}{11.30}\textbf{Gabarito: A.} O índice composto começa por \texttt{orgao_id}; filtrar por esse prefixo e depois intervalo em data tende a ser bom uso. Uso real depende do otimizador, estatísticas e SGBD.\end{solucao}
\begin{solucao}{11.31}\textbf{Gabarito: A.} Se valor total da nota é repetido em cada item, somá-lo na granularidade item multiplica o total. O erro é de compatibilidade entre medida e granularidade.\end{solucao}
\begin{solucao}{11.32}\textbf{Gabarito: A.} SCD tipo 2 preserva histórico criando nova versão da dimensão com validade/chave surrogate apropriada. Sobrescrever é comportamento típico de tipo 1.\end{solucao}
\begin{solucao}{11.33}\textbf{Gabarito: A.} Sob partição, CAP força trade-off entre consistência forte e disponibilidade para todas as requisições. As outras combinações não são o teorema.\end{solucao}
\begin{solucao}{11.34}\textbf{Gabarito: A.} Chave-valor é natural para lookup simples por chave. Se junções relacionais complexas são centrais, outro modelo pode ser mais apropriado.\end{solucao}
\begin{solucao}{11.35}\textbf{Gabarito: A.} Exactly-once depende do escopo da garantia. Efeitos externos podem exigir transação distribuída, outbox ou idempotência. A rede pode retransmitir e falhas continuam existindo.\end{solucao}
\begin{solucao}{11.36}\textbf{Gabarito: A.} Event time representa quando o evento ocorreu; watermarks e janelas tratam eventos tardios. Tempo de processamento pode distorcer análise temporal.\end{solucao}
\begin{solucao}{11.37}\textbf{Gabarito: A.} Padronização, validação e deduplicação com regra rastreável tratam consistência do identificador. Apagar sem evidência ou aceitar qualquer valor destrói qualidade.\end{solucao}
\begin{solucao}{11.38}\textbf{Gabarito: A.} Função sobre coluna pode impedir busca direta pelo índice, salvo suporte a índices funcionais ou reescrita otimizada. Não há garantia de index seek.\end{solucao}
\begin{solucao}{11.39}\textbf{Gabarito: A.} Toda mudança em dados pode exigir atualização de índices relacionados, aumentando I/O, CPU e espaço. Índices são trade-off, não aceleração gratuita.\end{solucao}
\begin{solucao}{11.40}\textbf{Gabarito: A.} O mesmo administrador pode alterar o objeto e destruir a evidência, fragilizando segregação e accountability. ETL, data lake e cardinalidade não tratam esse risco.\end{solucao}
\begin{solucao}{11.41}\textbf{Gabarito: A.} Retries são parte normal de sistemas resilientes; a solução é tornar efeito idempotente com chave estável, deduplicação e transação. Desabilitar retry reduz resiliência e não corrige causa.\end{solucao}
\begin{solucao}{11.42}\textbf{Gabarito: A.} Dado aberto sem metadados, periodicidade e versão perde contexto, comparabilidade e reprodutibilidade. Isso é problema de governança, não de infraestrutura.\end{solucao}
\begin{solucao}{11.43}\textbf{Gabarito: A.} Recarregar bilhões diariamente pode ser desnecessário. CDC/carga incremental e particionamento podem reduzir janela, desde que semântica e reconciliação sejam preservadas.\end{solucao}
\begin{solucao}{11.44}\textbf{Gabarito: A.} Sobrescrever atributo que precisa de histórico faz relatórios antigos mudarem retrospectivamente. O requisito pede estratégia de dimensão histórica, como SCD2 conforme caso.\end{solucao}
\begin{solucao}{11.45}\textbf{Gabarito: A.} Nome não é identificador confiável. Entity resolution deve usar atributos adicionais, score, regras de ambiguidade e validação para reduzir falsos positivos/negativos.\end{solucao}
\begin{solucao}{11.46}\textbf{Gabarito: A.} Ser ambiente interno não autoriza acesso irrestrito. Classificação, need-to-know, mínimo privilégio e trilha de acesso continuam necessários.\end{solucao}
\begin{solucao}{11.47}\textbf{Gabarito: A.} JOIN 1:N pode multiplicar medidas. O auditor deve conhecer granularidade, cardinalidade e reconciliar totais com fonte independente. DISTINCT pode mascarar problema sem resolver semântica.\end{solucao}
\begin{solucao}{11.48}\textbf{Gabarito: A.} PITR exige backup base, logs contínuos íntegros, retenção e procedimento testado. RAID ou mera existência de WAL não prova que a restauração funcionará.\end{solucao}
\begin{solucao}{11.49}\textbf{Gabarito: A.} Particionamento deve alinhar-se aos predicados e pruning. Muitas partições sem eliminação podem aumentar overhead. Partição não substitui índice.\end{solucao}
\begin{solucao}{11.50}\textbf{Gabarito: A.} Sem versão do código/parâmetros não há reprodutibilidade nem explicação confiável da mudança. Versionamento e lineage são controles centrais de engenharia de dados.\end{solucao}
